\item \textbf{Determinación de masa de cobre por activación de neutrones.}
Desea determinar la cantidad de átomos del isótopo $^{65}\text{Cu}$ en una pequeña muestra de material. Bombardea la muestra con neutrones hasta lograr que exactamente el $1\%$ de los núcleos de $^{65}\text{Cu}$ absorba un neutrón para transmutar en $^{66}\text{Cu}$ ($T_{1/2} = 5.10\text{ min}$). Al interrumpir el flujo, el $^{66}\text{Cu}$ decae; la mitad de ellos emite un rayo gamma característico de $1.04\text{ MeV}$. Después de $10.0\text{ min}$ de monitoreo, su detector registra una energía total de fotones gamma de $1.00 \times 10^4\text{ MeV}$.
\begin{enumerate}
    \item ¿Cuántos átomos de $^{65}\text{Cu}$ había originalmente en la muestra?
    \item Si la muestra contiene cobre natural ($30.8\%$ de $^{65}\text{Cu}$ y $69.2\%$ de $^{63}\text{Cu}$), estime la masa total de cobre metálico en la muestra.
\end{enumerate}